% บทนำภาษาไทย ใช้ใน main_short_th.tex (ฉบับอังกฤษอยู่ใน 01_introduction.tex); แก้ให้ตรงกันทั้งสองฉบับ
% ตัวเลขทั้งหมดมาจาก numbers.tex. ไม่ใช้ \emph เพราะฟอนต์ไทยไม่มีตัวเอียง.
\section{บทนำ}\label{sec:intro}

การสแกนห้องเป็นแบบจำลองสามมิติเป็นความสามารถที่มีอยู่แล้วใน iPhone~Pro และ iPad~Pro
(เช่น Apple RoomPlan~\cite{roomplan}) โดยอาศัยเซนเซอร์ LiDAR วัดความลึกเป็นหน่วยเมตร
แต่อุปกรณ์ที่มี LiDAR เป็นส่วนน้อย iPhone รุ่นที่ไม่ใช่ Pro และโทรศัพท์ Android แทบทุกรุ่นไม่มีเซนเซอร์นี้
ทางเลือกของอุปกรณ์เหล่านี้คือการประมาณความลึกจากภาพเดี่ยว (monocular depth estimation)
คือการทำนายความลึกจากภาพสีเพียงภาพเดียว ซึ่งพัฒนาอย่างรวดเร็วตั้งแต่ MiDaS~\cite{midas}
จนถึง Depth Anything~\cite{da1,da2} ทีมพัฒนาผลิตภัณฑ์จึงต้องรู้ว่า
\textbf{หากใช้โมเดลจากภาพเดี่ยวแทน LiDAR ห้องที่ได้จะคลาดเคลื่อนกี่เซนติเมตร และยังใช้งานแบบใดได้บ้าง}
เกณฑ์วัดมาตรฐานแบบรายภาพ~\cite{eigen14} ซึ่งเฉลี่ยความคลาดเคลื่อนของความลึกบนพิกเซลของภาพแต่ละภาพ
ตอบคำถามนี้ไม่ได้ และระบบสร้างแบบจำลองแบบครบวงจร (end-to-end)~\cite{atlas,neuralrecon,simplerecon}
ก็ตอบไม่ได้ เพราะออกแบบกระบวนการทั้งหมดรอบโมเดลของตนเอง

โมเดลรุ่นใหม่บางตัวให้ความลึกเป็นเมตรได้โดยตรง~\cite{da2,zoedepth,depthpro}
แต่ฝึกด้วยภาพที่ไม่ได้มาจากกล้องโทรศัพท์ เช่น \mbox{DA-v2} Metric-Indoor ฝึกด้วยภาพห้องสังเคราะห์จากชุดข้อมูล
Hypersim~\cite{hypersim} และในการทดลองของเรา ห้องที่สร้างด้วย \mbox{DA-v2} Metric-Indoor
คลาดเคลื่อนเฉลี่ย \rsMETRICCM{}~ซม. วิธีแก้โดยตรงคือการปรับจูน (fine-tune)
โมเดลด้วยความลึกจากอุปกรณ์เป้าหมาย คำถามคือข้อมูลสอน (label) นี้จะมาจากไหน
ระหว่างเครื่องสแกนเลเซอร์ระดับงานสำรวจ~\cite{faro} ซึ่งผู้พัฒนาแอปเพียงไม่กี่รายมีใช้
กับ LiDAR ที่มีอยู่แล้วในอุปกรณ์รุ่น Pro ซึ่งแอปเก็บความลึกได้ระหว่างการใช้งานปกติ
หากวิธีหลังได้ผล อุปกรณ์ที่มี LiDAR ก็จะเป็นครูสอนอุปกรณ์ที่ไม่มี LiDAR ได้

งานวิจัยนี้จึงตั้งคำถามว่า \mbox{DA-v2} Metric-Indoor ที่ปรับจูนด้วยความลึกจาก LiDAR ของ iPad
โดยไม่ใช้เครื่องสแกนเลเซอร์ในการฝึก แม่นยำกว่าโมเดลเดิมที่ไม่ได้ปรับจูนหรือไม่
เมื่อวัดเทียบกับเครื่องสแกนเลเซอร์ Faro ซึ่งเป็นค่าอ้างอิงของชุดข้อมูล ARKitScenes~\cite{arkitscenes}
บนห้องที่ไม่เคยใช้ฝึก โดยมีวัตถุประสงค์ 4 ข้อ ได้แก่
(1)~ปรับจูนโมเดลด้วยความลึกจาก LiDAR ของ iPad จำนวน 24 ห้อง ได้โมเดล \ftmain{}
แล้วเปรียบเทียบกับโมเดลเดิมบนห้องที่ไม่เคยใช้ฝึก โดยถือว่าดีกว่าเมื่อค่าเฉลี่ยความคลาดเคลื่อนต่ำกว่า
และผลต่างรายห้องมีนัยสำคัญทางสถิติที่ระดับ 5\,\% ตามการทดสอบ Wilcoxon signed-rank แบบ exact~\cite{wilcoxon}
(2)~วัดว่าแหล่งของข้อมูลสอนมีผลหรือไม่ โดยเปรียบเทียบข้อมูลสอนจาก Faro กับจาก LiDAR
บน 10 ห้องเดียวกัน (\ftfaro{} กับ \ftlidar{}) เปรียบเทียบข้อมูลสอนจาก LiDAR 10 ห้องกับ 24 ห้อง
และเปรียบเทียบกับโมเดลที่ใช้ความยาวโฟกัสของกล้องแทนข้อมูลสอนจากโดเมนเป้าหมาย (Depth Pro~\cite{depthpro})
(3)~ระบุตำแหน่งของ \ftmain{} เทียบกับแหล่งความลึกที่ใช้ได้ขณะใช้งาน ได้แก่ LiDAR ของ iPad
และโมเดลความลึกสัมพัทธ์ที่ปรับสเกลรายเฟรมด้วยจุดเบาบาง (sparse points) หรือปรับสเกลครั้งเดียวต่อห้อง
และ (4)~หาที่มาของความคลาดเคลื่อนที่การปรับจูนแก้ไม่ได้ โดยหาสเกลความลึกที่เหมาะกับแต่ละเฟรม
แล้วทดสอบว่าสเกลนั้นเปลี่ยนตามเวลา ซึ่งบ่งชี้การสะสมความคลาดเคลื่อน (drift) หรือเปลี่ยนตามสิ่งที่อยู่ในภาพ

ทุกการเปรียบเทียบวัดสองแบบ แบบแรกวัดเป็นแบบจำลองห้อง (mesh) ด้วย Chamfer distance~\cite{chamfer}
คือระยะเฉลี่ยจากแต่ละพื้นผิวถึงจุดที่ใกล้ที่สุดของอีกพื้นผิว และด้วย F-score ที่ 5~ซม.
(\fscore{})~\cite{tanks} ซึ่งรวมสัดส่วนของพื้นผิวที่สร้างได้ที่อยู่ภายใน 5~ซม. จากพื้นผิวอ้างอิง
กับสัดส่วนของพื้นผิวอ้างอิงที่ถูกสร้างขึ้นภายใน 5~ซม. แบบที่สองวัดรายเฟรมด้วย AbsRel~\cite{eigen14}
คือความคลาดเคลื่อนสัมพัทธ์เฉลี่ยของความลึก และด้วยสัดส่วนของพิกเซลที่ความลึกที่ทำนายต่างจากค่าจริงไม่เกิน 1.25 เท่า
(ไม่ว่าจะมากหรือน้อยกว่าค่าจริง) ซึ่งเรียกว่า \wone{}~\cite{eigen14}
สำหรับวัตถุประสงค์ข้อ (3) สัดส่วนของพื้นผิวอ้างอิงที่ถูกสร้างขึ้นภายใน 2, 5 และ 10~ซม.
ใช้แทนความแม่นยำที่งานวัดเพื่อปรับปรุงห้อง งานวางเฟอร์นิเจอร์ และงานดูภาพรวมของห้องต้องการ ตามลำดับ
และถือว่ารองรับงานนั้นได้เมื่อสัดส่วนนี้ถึง 90\,\%

ขอบเขตของงานวิจัยมีดังนี้ ข้อมูลเป็นห้องในบ้านจากชุดข้อมูล ARKitScenes ที่ถ่ายด้วย iPad~Pro รุ่นเดียว
(ภาพสีขนาด $640\times480$ พิกเซล และความลึกจาก LiDAR ขนาด $256\times192$ พิกเซล)
ค่าอ้างอิงคือความลึกจาก Faro ที่มากับชุดข้อมูล และไม่ใช้ LiDAR ของ iPad เป็นค่าอ้างอิง
การปรับจูนทำเฉพาะส่วน DPT head~\cite{dpt} ซึ่งเป็นตัวถอดรหัสความลึกของ \mbox{DA-v2} Metric-Indoor ขนาด Large
ด้วยสูตรการฝึกเดียว โดยตรึงส่วนเข้ารหัสภาพไว้ ประเมินผลบนห้องทดสอบ 6 ห้องทั้งแบบแบบจำลองห้องและแบบรายเฟรม
และอีก \rsVN{} ห้องแบบรายเฟรมเท่านั้น ทุกห้องไม่ซ้ำกับห้องที่ใช้ฝึกหรือใช้เลือก checkpoint
ทุกแหล่งความลึกผ่านกระบวนการแบบออฟไลน์เดียวกัน คือความลึกรายเฟรม การปรับสเกล
การรวมเป็นปริมาตร TSDF (truncated signed distance function)~\cite{curless96} และการสร้างพื้นผิว
แหล่งความลึกจึงเป็นตัวแปรเดียว นอกจากนี้ ทุกแหล่งและค่าอ้างอิงใช้ตำแหน่งกล้องชุดเดียวกัน
ซึ่งบันทึกโดย ARKit ความคลาดเคลื่อนของตำแหน่งกล้องจึงไม่อยู่ในผลการวัด สิ่งที่อยู่นอกขอบเขต ได้แก่
ระบบสร้างแบบจำลองที่เรียนรู้ทั้งกระบวนการ~\cite{atlas,neuralrecon,simplerecon}
กล้องโทรศัพท์รุ่นอื่นรวมถึง Android การประมวลผลแบบเรียลไทม์บนอุปกรณ์ และการประมาณตำแหน่งกล้อง

ประโยชน์ที่คาดว่าจะได้รับสอดคล้องกับวัตถุประสงค์ทีละข้อ ข้อแรก ผู้พัฒนาแอปจะรู้ว่าข้อมูลที่เก็บได้อยู่แล้ว
จากผู้ใช้อุปกรณ์รุ่น Pro ช่วยให้การสแกนห้องบนอุปกรณ์ที่ไม่มี LiDAR ดีขึ้นหรือไม่
โดยไม่ต้องปรับเทียบหรือใช้เซนเซอร์เพิ่มขณะใช้งาน ข้อสอง จะรู้ว่าการสร้างข้อมูลสอนจำเป็นต้องใช้เครื่องสแกนเลเซอร์หรือไม่
ข้อสาม จะรู้ว่าอุปกรณ์ที่ไม่มี LiDAR รองรับงานแบบใดได้บ้าง
พร้อมได้กระบวนการที่ประเมินแหล่งความลึกใหม่ได้ด้วยการแก้ค่าตั้งค่าเพียงรายการเดียว
และข้อสี่ จะรู้ว่าควรลงทุนพัฒนาส่วนใดต่อ ระหว่างสเกลรายเฟรม ข้อมูลฝึก หรือขนาดโมเดล
